\subsection{The full quadratic conductor at three axes}

\begin{statement}{P4.47}{Proved}
Consider an actual torus degeneration at a point where the
entire reduced critical curve is the three coordinate axes,
with generic transverse type $A_3$ on each axis. Suppose
the local Hessian--conductor defect is nine, the three
Hessian orders are two, and the quartic equation is the
square of $a xy+b xz+c yz$ with $abc\ne0$.
Then the normalization has only rational singularities
over that point. In particular it cannot contain the
contraction of a positive-genus horizontal section.
\end{statement}
\begin{proof}
Put $P=(xy,xz,yz)$ and $h=xyz$. The symbolic square is
$P^{(2)}=P^2+(h)$, and the actual conductor image ideal
$I$ contains $P^{(2)}$. Also $f\in P^{(2)}$.
Write $f=Ah+P^2$. On the $x$-axis, a nonzero constant
term of $A$ would give a leading rank-two transverse
Hessian, since the quadratic term is then a nonzero
multiple of $xyz$. Generic type $A_3$ has rank one.
Thus $A(0)=0$. Since $\mathfrak m h\subset P^2$,
we have $f\in P^2$.

The normalized conormal lattice is $tK$ on each axis,
so its projected colength is one. The curve has
$\delta=2$ and $R=0$. The stated defect therefore
forces the remaining normalization descent colength
to be one. The normalized quotient values form a
two-dimensional subspace of $\mathbb C^3$ with one
relation involving all three coordinates. Multiplying
a source generator that is a unit on one branch by
that branch's axis coordinate generates its positive
orders alone. Hence all positive-order branch functions
occur in the image, and the leading kernel vector
lifts to an actual element $g\in I$ with quadratic
part $a xy+b xz+c yz$.

We prove the exact ideal equality $I=(g,h)$.
On each coordinate plane, $g$ is the product of its
two coordinate functions times a unit. Thus
$I_0=(g,h)$ has radical $P$. Furthermore $P^2\subset I_0$:
modulo $h$, multiplication of $g$ by $xy$ gives a unit
times $(xy)^2$, and the other two squares follow in
the same way; all mixed products are multiples of $h$.
At the generic point of every axis, $I_0$ is the
tacnode conductor ideal, generated by a normal coordinate
and the square of the null coordinate. It agrees
generically with $I$. The ideal $I_0$ is a
codimension-two complete intersection, so
$\mathcal O/I_0$ is Cohen--Macaulay and $P/I_0$ is
torsionfree on the reduced curve. The map
$P/I_0\to P/I$ has finite-length kernel, which must
vanish. This proves $I=I_0$.

The actual conductor ring condition, with the hypotheses
verified in P4.41, gives $f\in I^2$ by the two-by-two
adjugate calculation of de Jong--van Straten,
Theorem~1.12. Write $f=A g^2+Bgh+Ch^2$.
Its quartic part makes $A(0)\ne0$. Completing the square
and absorbing $\sqrt A$ into the first factor gives
the exact equation
\[
                         f=G^2+h^2U.                 \tag{4.47a}
\]
The quadratic part of $G$ is nondegenerate, and
$(G,h)=(g,h)$. If $U(0)\ne0$, the two factors
$G\pm h\sqrt{-U}$ have nondegenerate quadratic
Hessians. Each is an isolated $A_1$ surface singularity,
so the full normalization has rational singularities.

We exclude $U(0)=0$ by a deformation that keeps the
three axes fixed. Add an arbitrarily small generic
linear form to $U$, and then replace $U$ by
$U-\epsilon$. On a fixed boundary the transverse
$A_3$ coefficient is a unit, so these changes retain
the boundary transverse type. A diagonal linear change
puts the quadratic part of $G$ in the form
$G_2=xy+xz+yz$, changing $h$ only by a nonzero constant.
Let $L$ be the resulting linear part of $U$.
Under $(x,y,z)=\epsilon(X,Y,Z)$, equation (4.47a) becomes
\[
 \frac{f_\epsilon}{\epsilon^4}
   =G_\epsilon(X,Y,Z)^2+\epsilon^3(XYZ)^2
                          \bigl(L(X,Y,Z)-1+O(\epsilon)\bigr),
 \qquad G_\epsilon=G_2+O(\epsilon).
\]
Near a nonzero point of the smooth cone $G_2=0$,
use $G_\epsilon$ as the normal coordinate. The
parameter Morse centre is $O(\epsilon^3)$.
Substitution leaves the tangential function
$\epsilon^3(\Phi+O(\epsilon))$, where
\[
                 \Phi=(XYZ)^2(L-1)|_{G_2=0}.
\]

We verify a nonempty open set of choices of $L$ for
which $\Phi$ has a nondegenerate critical point of
nonzero value. Parametrize the projective conic by
\[
              [X:Y:Z]=[t(t+1):t+1:-t].
\]
For $L=X+Y+Z$, set
$H(t)=-t^2(t+1)^2$ and $\ell(t)=t^2+t+1$.
In radial coordinate $r$,
$\Phi=r^6H(t)^2(r\ell(t)-1)$.
Radial criticality gives $r=6/(7\ell(t))$, and
\[
 \frac{d}{dt}\left(\frac{H}{\ell^3}\right)
  =\frac{t(t-1)(t+1)(t+2)(2t+1)}{(t^2+t+1)^4}.
\]
The three critical points $t=1,-2,-1/2$ are simple
and avoid zeros and poles. Their radial coordinates
are $2/7,2/7,8/7$. At each point,
\[
       \Phi=-\frac{1024}{823543},\qquad
       \det\operatorname{Hess}_{r,t}\Phi
                     =\frac{4718592}{1977326743}.
\]
These numbers are nonzero. Simplicity and avoidance
of zeros and poles are nonvanishing algebraic conditions
on the coefficients of $L$. The explicit example proves
that they hold on a nonempty Zariski open set, so an
arbitrarily small generic linear perturbation of any
initial linear part is sufficient.

After eliminating the normal Morse coordinate, divide
the two tangential critical equations by $\epsilon^3$.
They are analytic in $\epsilon$ and have invertible
Jacobian at the displayed points. The implicit-function
theorem lifts them to nondegenerate three-variable
critical points of $f_\epsilon$. Their critical values
have leading term $\epsilon^7$ times the displayed
nonzero value. They tend to the original point.
Since the deformation retains boundary $A_3$, the
nested boundary transport embeds a Morse vanishing
sphere of square $-2$ into a nearby global torus fibre.
This contradicts asphericity. Therefore $U(0)=0$
cannot occur, proving the assertion.
\end{proof}

The conic derivative, all three exact critical values
and Hessian determinants, and the related conductor
matrices are executed in
\texttt{code/part4/a3\_three\_axes.py}.
The output reports PASS. The working record explicitly
rejects the earlier proposed change $G\mapsto G-\epsilon$:
with $h$ fixed that change replaces the boundary tacnode
by two ordinary double points and does not supply the
required boundary transport.
